% =====================================================================
%  Guide — Environment Setup: ECS + SSH Proxy + MaaS Gateway + Huawei Project
%  Compile with XeLaTeX:
%      xelatex setup-guide.tex   (run 2x for the TOC)
% =====================================================================
\documentclass[english]{guide}

% --- Document configuration -----------------------------------------
\setguidetitle{Guide: Environment Setup for Huawei Cloud Development}
\setheadertitle{Huawei Cloud -- ECS + SSH + MaaS Gateway Setup}
\setcovertext{Huawei Technologies CO., LTD}
\setdocversion{2.0.0}

\begin{document}

% ---- COVER ---------------------------------------------------------
\makecover

% ---- TABLE OF CONTENTS ---------------------------------------------
\sumario

% ---- BODY (restarts page numbering at 1) ---------------------------
\startbody

% =====================================================================
%  Chapter 1 — Provision an ECS Instance
% =====================================================================
\section{Provision an ECS Instance on Huawei Cloud}

\begin{objetivos}
  \objgeral{Provision an Elastic Cloud Server (ECS) on Huawei Cloud that will
  host a complete coding stack environment: the
  \weblink{https://github.com/wallacelw/oh-my-coding-maas-gateway}{oh-my-coding-maas-gateway}
  (a LiteLLM-based proxy with load balancing, virtual keys, and Grafana
  observability), governance tooling, and multiple coding agents --- with
  OpenCode as the primary one. By the end of this chapter you will have a
  running ECS instance with a public IP accessible over SSH.}

  \prerequisitos
  \begin{itemize}
    \item Active Huawei Cloud account with IAM credentials.
    \item Permission to create ECS instances, EIPs, and security groups.
    \item A key pair created in the Console, or permission to create one.
    \item A web browser (Chrome, Firefox, or Edge).
  \end{itemize}
\end{objetivos}

\begin{infobox}
This guide assumes the \code{sa-brazil-1} region. Adjust the region if your
account is registered elsewhere.
\end{infobox}

\subsection{Create the ECS instance}

\objpratica{Launch a single ECS instance using the Console wizard.}

\passoapasso
\begin{enumerate}
  \item Log in to the \weblink{https://console.huaweicloud.com/}{Huawei Cloud
        Console} and select the \code{sa-brazil-1} region in the top-right
        corner.

  \item Open the service catalog and select \textbf{Elastic Cloud Server
        (ECS)}.

  \item Click \textbf{Create ECS} and fill in the basic configuration:

  \begin{codigo}
Name:           dev-server-01
Flavor:         s6.large.1 (2 vCPU, 2 GB RAM)
Image:          Ubuntu 22.04 Server 64bit
Disk:           40 GB General Purpose SSD
VPC:            vpc-default
Subnet:         subnet-default
Security Group: sg-ssh (allows inbound TCP 22)
  \end{codigo}

  \item Under \textbf{Login Configuration}, select the key pair you created
        earlier. Click \textbf{Create Now} and wait for the status to change
        to \textbf{Running}.

  \badge{Important}

  \begin{aviso}
  The instance is assigned a public IP (EIP). If you need a fixed IP,
  allocate a dedicated EIP and bind it after creation. The EIP is released
  when the instance is deleted.
  \end{aviso}
\end{enumerate}

\subsection{Configure the security group}

If you do not already have a security group that allows SSH access:

\passoapasso
\begin{enumerate}
  \item In the Console, go to \textbf{VPC} $\rightarrow$ \textbf{Security
        Groups} and click \textbf{Create Security Group}.
  \item Name it \code{sg-ssh} and add an inbound rule:

  \begin{codigo}
Protocol:    TCP
Port:        22
Source:      0.0.0.0/0
Description: Allow SSH access
  \end{codigo}

  \item Associate the security group with your ECS instance.
\end{enumerate}

\begin{aviso}
Opening SSH to \code{0.0.0.0/0} allows access from any IP. For production,
restrict the source to your specific IP address.
\end{aviso}

\subsection{Verify the instance}

On the ECS details page, copy the \textbf{EIP} field. This is the address you
will use to connect in the next chapter.

% =====================================================================
%  Chapter 2 — Connect via SSH with Proxy
% =====================================================================
\section{Connect via SSH with Proxy for VS Code}

\begin{objetivos}
  \objgeral{Establish an SSH connection to the ECS instance through a proxy
  (if required by your network) and configure VS Code Remote-SSH for
  remote development.}

  \prerequisitos
  \begin{itemize}
    \item Chapter 1 completed — ECS instance running with a public IP.
    \item The private key file (\code{.pem}) corresponding to the key pair.
    \item VS Code installed locally with the Remote-SSH extension.
    \item If behind a corporate proxy: the proxy host, port, and credentials.
  \end{itemize}
\end{objetivos}

\subsection{Configure SSH with a proxy}

If your network requires a proxy to reach the internet, configure SSH to
tunnel through it using \code{ProxyCommand}.

\passoapasso
\begin{enumerate}
  \item Open or create the SSH config file:

  \begin{codigo}[bash]
nano ~/.ssh/config
  \end{codigo}

  \item Add the following configuration (adjust the proxy and EIP values):

  \begin{codigo}
Host huawei-dev
    HostName <EIP>
    User ubuntu
    IdentityFile ~/.ssh/huawei_key.pem
    ProxyCommand nc -X connect -x proxy.corp.com:8080 %h %p
    ServerAliveInterval 60
    ServerAliveCountMax 3
  \end{codigo}

  \begin{infobox}
  Replace \code{proxy.corp.com:8080} with your actual proxy host and port.
  If your proxy requires authentication, use:
  \code{ProxyCommand nc -X connect -x user:pass@proxy.corp.com:8080 \%h \%p}
  \end{infobox}

  \item If you do \textbf{not} need a proxy, use this simpler config:

  \begin{codigo}
Host huawei-dev
    HostName <EIP>
    User ubuntu
    IdentityFile ~/.ssh/huawei_key.pem
    ServerAliveInterval 60
    ServerAliveCountMax 3
  \end{codigo}

  \item Set the correct permissions on the key file and test the connection:

  \begin{codigo}[bash]
chmod 600 ~/.ssh/huawei_key.pem
ssh huawei-dev
  \end{codigo}
\end{enumerate}

\begin{dica}
The \code{ServerAliveInterval} and \code{ServerAliveCountMax} settings keep
the connection alive through proxies that drop idle connections.
\end{dica}

\subsection{Connect VS Code via Remote-SSH}

\passoapasso
\begin{enumerate}
  \item Install the \textbf{Remote - SSH} extension in VS Code
        (publisher: Microsoft).

  \begin{codigo}[bash]
code --install-extension ms-vscode-remote.remote-ssh
  \end{codigo}

  \item Press \code{F1}, type \textbf{Remote-SSH: Connect to Host...}, and
        select \code{huawei-dev} from the list.

  \item VS Code opens a new window connected to the remote instance. Click
        \textbf{Open Folder} and select \code{/home/ubuntu} (or your working
        directory).

  \item Verify the connection by opening a terminal in VS Code
        (\code{Ctrl+`}) and running:

  \begin{codigo}[bash]
uname -a
df -h
  \end{codigo}
\end{enumerate}

% =====================================================================
%  Chapter 3 — Install oh-my-coding-maas-gateway
% =====================================================================
\section{Install the MaaS Coding Gateway}

\begin{objetivos}
  \objgeral{Install the oh-my-coding-maas-gateway on the ECS instance.
  This is a self-hosted LiteLLM proxy that routes Huawei ModelArts MaaS
  models to coding tools (opencode, Codex CLI, Claude Code, Pi agent) with
  load balancing, virtual keys, and Prometheus + Grafana observability.}

  \prerequisitos
  \begin{itemize}
    \item Chapter 2 completed — VS Code connected to the ECS via Remote-SSH.
    \item Docker and Docker Compose available (the installer handles this).
    \item A Huawei ModelArts MaaS API key.
  \end{itemize}
\end{objetivos}

\begin{infobox}
The gateway repository is at
\weblink{https://github.com/wallacelw/oh-my-coding-maas-gateway}{wallacelw/oh-my-coding-maas-gateway}.
It supports 4 models: \code{glm-5.2}, \code{glm-5.1},
\code{deepseek-v4-pro}, and \code{deepseek-v4-flash}.
\end{infobox}

\subsection{Run the bootstrap installer}

\objpratica{Install the gateway using the one-command bootstrap script.}

\passoapasso
\begin{enumerate}
  \item In the VS Code remote terminal, run the bootstrap script:

  \begin{codigo}[bash]
curl -fsSL https://raw.githubusercontent.com/wallacelw/oh-my-coding-maas-gateway/main/scripts/bootstrap.sh | bash
  \end{codigo}

  \item The installer will:
    \begin{itemize}
      \item Prompt for your Huawei MaaS API key.
      \item Show an interactive menu to select which coding tools to install
            (opencode, Codex CLI, Claude Code, Pi agent).
      \item Auto-install prerequisites (Docker, Node.js, etc.).
      \item Deploy the proxy + observability stack (LiteLLM, PostgreSQL,
            Prometheus, Grafana).
      \item Configure each tool with its own virtual key.
      \item Run validation and offer to install the companion skill.
    \end{itemize}

  \badge{Important}

  \begin{aviso}
  Estimated time: \textasciitilde5 min for a fresh install, \textasciitilde2 min
  for an upgrade. If you are behind a proxy, ensure \code{http\_proxy} and
  \code{https\_proxy} environment variables are set before running the script.
  \end{aviso}
\end{enumerate}

\subsection{Verify the installation}

\passoapasso
\begin{enumerate}
  \item Check that the proxy is running:

  \begin{codigo}[bash]
curl http://localhost:4000/health
  \end{codigo}

  \item Open the Grafana dashboard at \code{http://localhost:3000} to
        monitor usage and metrics.

  \item Open the LiteLLM Admin UI at \code{http://localhost:4000/ui} to
        manage virtual keys and model routing.

  \item Launch your coding tool:

  \begin{codigo}[bash]
opencode          # or:  codex  or:  claude --bare  or:  pi
  \end{codigo}
\end{enumerate}

\begin{dica}
To upgrade the gateway later, re-run the same bootstrap command. It detects
existing installs, compares versions, and pulls updates — all secrets and
data are preserved.
\end{dica}

% =====================================================================
%  Chapter 4 — Install the Huawei Project
% =====================================================================
\section{Install the Huawei Document Templates Project}

\begin{objetivos}
  \objgeral{Clone the Huawei document templates repository on the ECS
  instance, run the setup script, and verify that the LaTeX toolchain and
  the coding agent skill are ready for use.}

  \prerequisitos
  \begin{itemize}
    \item Chapter 3 completed — MaaS gateway installed and running.
    \item Git installed (\code{sudo apt install git}).
    \item The repository URL for the Huawei document templates project.
  \end{itemize}
\end{objetivos}

\subsection{Clone and set up the project}

\passoapasso
\begin{enumerate}
  \item Clone the repository:

  \begin{codigo}[bash]
cd ~
git clone <repo-url> Huawei-doc-templates
cd Huawei-doc-templates
  \end{codigo}

  \item Run the setup script (installs XeLaTeX, latexmk, fonts, the coding
        agent skill, and the VS Code LaTeX Workshop extension):

  \begin{codigo}[bash]
./install.sh
  \end{codigo}

  \begin{aviso}
  The script runs \code{sudo apt-get install} for TeX Live packages. If you
  do not have sudo access, ask your administrator to install the packages
  listed in \code{install.sh}.
  \end{aviso}

  \item Verify the LaTeX toolchain:

  \begin{codigo}[bash]
xelatex --version
latexmk --version
  \end{codigo}
\end{enumerate}

\subsection{Compile the sample guides}

\passoapasso
\begin{enumerate}
  \item Compile the Portuguese sample:

  \begin{codigo}[bash]
cd templates/guide/samples/pt
latexmk main.tex
  \end{codigo}

  \item Compile the English sample:

  \begin{codigo}[bash]
cd ../en
latexmk main.tex
  \end{codigo}

  \item Verify the PDFs were generated:

  \begin{codigo}[bash]
ls -l templates/guide/samples/pt/main.pdf
ls -l templates/guide/samples/en/main.pdf
  \end{codigo}
\end{enumerate}

\subsection{Create your first guide with the skill}

\passoapasso
\begin{enumerate}
  \item In the VS Code remote terminal, launch your coding agent:

  \begin{codigo}[bash]
cd ~/Huawei-doc-templates
opencode          # or: codex, claude --bare, pi
  \end{codigo}

  \item Run the guide skill:

  \begin{codigo}
/skill huawei-template-guide
  \end{codigo}

  \item The skill will ask for:
    \begin{itemize}
      \item \textbf{Title} — e.g. ``Provisioning an OBS Bucket''
      \item \textbf{Language} — Portuguese or English
      \item \textbf{Project name} — used as the folder name
      \item \textbf{Location} — where to create the project folder
    \end{itemize}

  \item The skill generates the \code{.tex} file, creates a
        \code{.latexmkrc} with the correct \code{TEXINPUTS}, compiles the
        document, and reports the page count.

  \badge{Ready}
\end{enumerate}

\begin{infobox}
The skill is auto-discovered via \code{opencode.json} which registers
\code{templates/} as a skill discovery path. No manual skill installation
is needed when working inside the repo.
\end{infobox}

\subsection{Write content using template commands}

The template provides commands for common guide elements:

Headings (automatic numbering, H1 starts on a new page):

\begin{codigo}
\section{Chapter Title}          % H1: giant number + title + rule
\subsection{Section Title}        % H2
\subsubsection{Subsection Title}  % H3
\end{codigo}

Objectives block (closes with a horizontal rule):

\begin{codigo}
\begin{objetivos}
  \objgeral{...}
  \prerequisitos
  \begin{itemize}
    \item ...
  \end{itemize}
\end{objetivos}
\end{codigo}

Callout boxes and inline commands:

\begin{codigo}
\begin{aviso} ... \end{aviso}     % amber warning box
\begin{dica}  ... \end{dica}      % green tip box
\begin{infobox} ... \end{infobox} % blue info box
\end{codigo}

\begin{codigo}
\code{inline code}                % inline monospace
\imagem{assets/screenshot.png}    % centered image
\nota{Italic observation.}        % note paragraph
\weblink{https://...}{link text}  % blue clickable link
\badge{New}                       % inline red label
\end{codigo}

% =====================================================================
%  Chapter 5 — Clean Up
% =====================================================================
\section{Clean Up}

\begin{objetivos}
  \objgeral{Remove the resources created in this guide to avoid ongoing
  charges on your Huawei Cloud account.}
\end{objetivos}

\passoapasso
\begin{enumerate}
  \item In the Console, navigate to \textbf{Elastic Cloud Server}.
  \item Select the instance \code{dev-server-01}, click \textbf{More} and
        choose \textbf{Delete}.
  \item Confirm deletion. Optionally release the EIP if it was dedicated.
  \item Remove the SSH config entry from \code{\textasciitilde/.ssh/config}.
\end{enumerate}

\begin{aviso}
If you used Terraform to provision the infrastructure, run
\code{terraform destroy} instead to remove all resources declared in the
configuration.
\end{aviso}

\nota{This guide was generated using the \code{guide.cls} template. For the
full command reference, see the template's README.md or run
\code{/skill huawei-template-guide} in your coding agent.}

% =====================================================================
%  Changelog
% =====================================================================
\section{Changelog}

\begin{changelog}
  \changelogentry{2.0.0}{2026-08-05}{
    \item Replaced OpenCode with oh-my-coding-maas-gateway as the default
          coding environment.
    \item Added changelog and versioning support to the template
          (\code{\textbackslash setdocversion}, \code{\textbackslash setdocdate},
          \code{changelog} environment).
    \item Renamed \code{info} environment to \code{infobox} to avoid
          package name collisions.
    \item Fixed \code{\textbackslash codigoarquivo} implementation for
          external file code blocks.
    \item Added GitHub Actions CI to compile samples on every push.
  }
  \changelogentry{1.0.0}{2026-08-05}{
    \item Initial version.
    \item Added ECS provisioning on Huawei Cloud.
    \item Added SSH with proxy configuration for VS Code Remote-SSH.
    \item Replaced OpenCode with oh-my-coding-maas-gateway (LiteLLM proxy
          with load balancing, virtual keys, and Grafana observability).
    \item Added Huawei document templates project setup.
    \item Added changelog and versioning support to the template.
  }
\end{changelog}

\end{document}
